\chapter{Outlook}
\label{ch:outlook}

\section{One principle, many faces}

We have now followed a single idea through six concrete settings, and it is worth
gathering the threads before asking what lies beyond them. In every case the
network grew by a \emph{local} rule --- a newcomer acting only on a node it had
reached and that node's immediate neighbors --- and in every case the same family
of global regularities emerged unbidden.

Searchable networks, such as the citation graph and the World Wide Web, grow by
exploration: a newcomer surfs along links and connects to what it finds. Modeling
the surfer as the very random walk that defines Google's PageRank, we found that a
page's visiting probability rises in a straight line with its in-degree, so that
the local act of walking manufactures an effective ``rich get richer.'' From it
follow a power-law degree distribution whose exponent is fixed purely by how the
network is explored, a clustering hierarchy in which cliquishness falls with
degree, and generically disassortative correlations
(Chapter~\ref{ch:search}).

Social networks grow by introduction: a friend introduces a friend, closing a
triangle. Formalizing the friend-of-a-friend as a \emph{potential link}, born in
proportion to degree, we recovered the same rich-get-richer and the same
clustering hierarchy --- but now with \emph{assortative} correlations, the
well-connected keeping one another's company (Chapter~\ref{ch:friends}). The
contrast with the searchable networks is that chapter's quiet lesson: the sign of
the degree correlations is not a free parameter but a readout of whether the rule
reaches outward, to the periphery, or inward, within a neighborhood.

Biological networks grow by copying: a gene is duplicated, its protein inheriting
all of the ancestor's interactions, and the copies then diverge. Because a node
gains a link whenever one of its neighbors is duplicated, copying too is a
rich-get-richer engine; self-interactions seed the clustering hierarchy, and
coupling divergence to duplication adds a density transition and a multifractal
degree distribution born of inheritance, alongside the disassortative mixing
measured in real protein networks (Chapter~\ref{ch:duplication}). Here the three
mechanisms met in one sentence: a rule that attaches to a \emph{neighbor} of a
randomly chosen node attaches in proportion to degree, because a well-connected
node simply belongs to more neighborhoods. Preferential attachment, the
clustering hierarchy, and the degree correlations are all the generic harvest of
locality.

The same copying, carried into project schedules, refines each activity either by
duplicating it into parallel copies or by splitting it into specialized
sequential parts. The rule predicts identical power laws for an activity's
predecessors and successors, a critical path that vanishes as a fraction of large
projects, and a fractal rather than small-world geometry --- all borne out in
seventy-seven real construction projects --- and it relocates the risk of delay
from the celebrated critical path to a percolation threshold set by the whole
network's average degree (Chapter~\ref{ch:projects}).

Finally, locality proved to shape structure above the level of the individual
node. The same rules that set the degrees also segregate the network into
communities --- even when every node is identical --- a fact we made precise
through the Ramsey community number, the size beyond which communities are all but
certain (Chapter~\ref{ch:communities}); and they thread the network with redundant
shortest paths, multiplying the routes between nodes far faster than chance and in
lockstep with those very communities (Chapter~\ref{ch:distance}).

And with that, the book comes full circle. We set out asking where the global
regularities of real networks come from --- their power-law degrees, their
small-world distances, their clustering hierarchy, their degree correlations ---
and answered, one mechanism at a time, that none of them need be imposed from
above: each emerges from a way of growing in which no node ever uses more than
local knowledge. A node that surfs to a neighbor, befriends a friend, is copied
from an ancestor, or refines a task asks nothing of the network as a whole --- and
yet the networks these humble acts build carry, in the aggregate, the full and
unmistakable signature of real complex systems.

\section{What remains to be done}

If the thesis of this book is right, it opens more questions than it closes.

Several are concrete. The exponential growth of shortest path multiplicity under
the duplication--split rule --- so much faster than the log-quadratic growth of
every other rule we studied --- still wants a full explanation; the
square-counting argument shows why copying outruns triadic closure, but not why it
runs away exponentially. The question is at least better posed than it was, for we
now know what the ceiling is: on the diamond lattice, where the branching can be
solved exactly, the multiplicity grows as a stretched exponential in $\sqrt{n}$,
and every stochastic rule falls below that bound (Chapter~\ref{ch:distance}). What
remains is to explain where between the log-quadratic floor and the deterministic
ceiling a given rule lands, and why copying alone climbs so far.

A second concrete gap sits beside it. The mean-field theory of the multiplicity
knows the network only through its excess degree $c_2$, and that is demonstrably
not enough: real networks fall on either side of the $c_2/e$ line according to the
sign of their degree correlations, the assortative co-authorship graph above it and
the disassortative Internet below. A theory that carries the correlations as well
as the degrees would close the gap between the calculation and the data, and it
would say something the models cannot yet say about which real systems are
route-rich and which are not.

The Ramsey community number has moved further still, though not in the direction
one might have hoped. That it depends on the method used to detect the communities
was always plain. Now that it can be computed exactly, we can say how far that
dependence goes: on one and the same web, under one and the same partition, a
plain block model puts $r_\kappa$ at $1095$ and a degree-corrected one at $42$, and
the two disagree about which cut is best (Chapter~\ref{ch:communities}). The dependence is not a rounding error to be
argued away; it is a factor of twenty-five, and it means a bare Ramsey number
quoted without its detection rule is not a property of the network at all. Yet the
same calculations hand us the beginnings of a repair. Underneath the
method-dependent threshold there is a method-independent quantity: the evidence
density flows to $\ln K$ per link, the entropy of resolving which of $K$
communities a link belongs to, and that value knows nothing of the lattice, the
degrees, or the priors. A definition of emergence built on the invariant rather
than on the crossing --- on the rate at which evidence accumulates rather than on
the size at which it clears an arbitrary bar --- would put the whole idea on firmer
ground.

One question spans both chapters at once, and it may be the most inviting of all.
The ring --- the barest case we can solve, a lattice with no growth and no
randomness in it --- turns out to be split into communities exactly when its links
outnumber its nodes, which is to say exactly when it carries an extensive number of
independent cycles. And cycles are also what multiply its shortest paths. Whether
that accounting is general, or an accident of the ring's simplicity --- whether the
surplus of links over nodes is the common currency of both communities and
redundant routes --- would fold the last two chapters into a single statement. The
book's thesis makes the question natural; it does not answer it.

Two smaller loose ends deserve naming, since both have the flavor of a clue rather
than a nuisance. The number of communities that best describes a self-similar
network appears to grow as $\sqrt{n}$, the geometric midpoint of the network's own
hierarchy --- but that is a numerical observation over the sizes we can reach, not
a theorem, and a proof would be worth having, not least because the same
$\sqrt{n}$ governs the shortest paths on the same lattice. And the Ramsey number's
climb with cycle length carries a persistent even--odd alternation, odd cycles
bringing communities on sooner than the even cycles around them. It is too regular
to be noise and has no explanation.

Even the self-organized critical state of the exhaustive searcher, where a
scale-free network appears across a whole range of the control parameter rather
than at one fine-tuned value, deserves to be connected more tightly to the wider
physics of criticality.

Others are programmatic, and the largest of them is the price we paid for solving
anything exactly at all. Every closed form in the last two chapters --- the ring's
evidence, the diamond's renormalization group, the pseudofractal's counts --- lives
on a \emph{deterministic} lattice, a fixed construction rather than a grown
network, with the candidate partition handed to the calculation rather than found
by an algorithm. That is what made the algebra possible, and it is also the
catch: those results bound and illuminate the simulated ones, but they do not
reproduce them. Whether the renormalization-group map survives off the lattice ---
exactly, or on average, for networks grown at random by a local rule --- sets how
far this reading of emergence extends. It is a demanding question, because on
disordered networks detectability is itself a genuine phase transition, and the
threshold we have been computing would become the edge of a phase rather than the
crossing of a bar. The reward would be a theory of emergence for the networks we
actually care about, rather than for their crystalline cousins.

Throughout, we read each local rule \emph{forward}, from
mechanism to statistics; the empirical chapters hint at the reverse problem ---
reading the rule \emph{backward}, inferring from a single observed network which
local moves grew it, and with what parameters, as we did when fitting a
duplication index to project schedules. A general inverse theory, turning measured
correlations, clustering hierarchies and community structure into an estimate of
the underlying rule, would make these models genuinely diagnostic. And real
networks are surely not grown by one rule in isolation: a social graph is part
introduction and part search, a cellular network part duplication and part fresh
interaction. How competing local moves combine --- whether their signatures simply
add, or interfere to produce something neither would alone --- is largely open.

The setting, too, can be widened. The networks of this book are mostly
unweighted, undirected and static, whereas real ones carry weights, directions,
timestamps, and several kinds of link at once. Whether the locality principle
survives into weighted, temporal and multilayer networks --- and what new
regularities it might generate there --- is a natural frontier. So is the coupling
of growth to \emph{function}: the percolation of delays through a project schedule
was a first glimpse of a dynamics running on a network that local rules had built,
and the interplay between how a network grows and what it then does promises more
than either studied alone.

None of these questions threatens the central claim; each extends it. Local rules
are the driving mechanism; everything else follows.
